\chapter{Background and Literature Review}
\label{ch:related-work}

This chapter provides the foundational background and surveys the literature relevant to this thesis. We organize the review into two main sections: \S\ref{sec:rw-concepts} introduces the fundamental concepts and formulations along the video retrieval pipeline, while \S\ref{sec:rw-literature} surveys existing work and identifies research gaps.

\section{Basic Concepts}
\label{sec:rw-concepts}

This section introduces the fundamental concepts and formulations underlying this thesis, organized along the core pipeline of video retrieval: from problem formulation (\S\ref{sec:bc-ir}--\S\ref{sec:bc-vr}), to the embedding models that produce representations (\S\ref{sec:bc-embedder}), to the spatiotemporal preprocessing required before videos enter the model (\S\ref{sec:bc-spatiotemporal}), and finally to the large-scale retrieval systems that deploy these components in practice (\S\ref{sec:bc-anns}--\S\ref{sec:bc-industrial}).

\subsection{Information Retrieval}
\label{sec:bc-ir}

Information retrieval (IR) is the task of identifying items from a collection that satisfy a user's information need. Formally, given a query $q$ and a corpus $\mathcal{C} = \{c_1, c_2, \ldots, c_N\}$ of $N$ candidate items, the goal is to retrieve the subset of items most relevant to $q$ according to some relevance criterion. In the ranking formulation, the system produces an ordered list of items sorted by their estimated relevance scores $s(q, c_i)$.

The effectiveness of an IR system depends on two core components: a \textbf{representation function} that encodes queries and items into a comparable form, and a \textbf{similarity function} that measures the relevance between their representations. The choice of representation determines what aspects of the data can be captured, while the similarity function determines how these representations are compared.

\subsection{Video Retrieval}
\label{sec:bc-vr}

Video retrieval extends the general IR framework to the domain of video content, where the corpus consists of videos and queries can take multiple modalities. Given a query $q$ and a video corpus $\mathcal{V} = \{v_1, v_2, \ldots, v_N\}$, the task is to identify videos whose content is relevant to $q$~\citep{zhu2023deep}.

\subsubsection{Query Modalities}

Video retrieval systems support diverse query formats:

\textbf{Textual Video Retrieval} (a.k.a., Text-to-Video Retrieval). The query $q$ is a natural language description, and the system retrieves videos matching the textual description. This is the most widely studied formulation, with benchmarks such as MSR-VTT~\citep{xu2016msr} and DiDeMo~\citep{anne2017localizing} establishing the standard evaluation paradigm.

\textbf{Composed Video Retrieval}. The query $q$ combines multiple modalities, typically a text description and a reference image or video. This formulation supports scenarios where users provide both visual examples and textual specifications to refine their search intent~\citep{ventura2024covr,hummel2024egocvr}.

\textbf{Visual Video Retrieval}. The query $q$ is a pure visual input (image or video), and the system retrieves semantically similar videos from the corpus~\citep{liu2021activity}.

\subsubsection{Semantic Domains}

Beyond query modality, video retrieval tasks vary in the granularity and nature of the semantic match required:

\textbf{Coarse-grained retrieval} involves matching high-level semantics, such as overall video themes or general content categories. Short-clip benchmarks like MSR-VTT typically fall into this category.

\textbf{Fine-grained retrieval} requires understanding detailed spatial relationships or temporal dynamics within videos, such as identifying specific objects, actions, or camera movements~\citep{xu2025carebench,lin2025towards}.

\textbf{Partially relevant retrieval} addresses scenarios where only a portion of a video is relevant to the query~\citep{dong2022partially}. Benchmarks such as TVR~\citep{lei2020tvr} and ActivityNet Captions~\citep{krishna2017dense} provide fine-grained temporal annotations for moment localization.

\textbf{Long-video retrieval} extends the temporal scope to hour-long videos, where relevant content may occupy only a small fraction of the total duration~\citep{cai2025lovr}. This setting introduces challenges of temporal sparsity and multi-granular query understanding.

\subsection{Video Embedding Models}
\label{sec:bc-embedder}

Modern video retrieval systems rely on embedding-based approaches, where queries and videos are mapped to dense vector representations in a shared semantic space, reducing retrieval to nearest neighbor search in the embedding space.

\subsubsection{Embedding Formulation}

Let $E_q: \mathcal{Q} \rightarrow \mathbb{R}^d$ and $E_v: \mathcal{V} \rightarrow \mathbb{R}^d$ denote the query and video embedding functions, which map inputs to $d$-dimensional vectors. The relevance between a query $q$ and a video $v$ is computed as the cosine similarity between their embeddings:
\begin{equation}
    s(q, v) = \frac{E_q(q) \cdot E_v(v)}{\|E_q(q)\| \, \|E_v(v)\|}
\end{equation}

The embedding functions are implemented as neural network encoders. Depending on the underlying architecture, these encoders can be broadly categorized into two families.

\subsubsection{CLIP-based Embedding Models}

The first family builds upon CLIP~\citep{radford2021learning}, a vision-language model pretrained via contrastive learning on large-scale image-text pairs. CLIP employs a \textbf{dual-encoder architecture}: a text encoder and an image encoder project their respective inputs into a shared embedding space where cross-modal similarity is measured by cosine similarity.

To extend CLIP to video retrieval, researchers have introduced temporal modeling into the visual encoder. These CLIP-based models retain the dual-encoder paradigm: a \textbf{text encoder} processes natural language queries, while a separate \textbf{video encoder} processes sampled frames, and the two encoders are trained to align text and video representations through contrastive objectives. A characteristic of CLIP-based models is that the query and video encoders are architecturally distinct modules tailored to their respective modalities. While effective for text-to-video retrieval, this design limits their ability to handle composed queries or visual queries, as the text encoder cannot natively process visual inputs and vice versa.

\subsubsection{MLLM-based Embedding Models}

The second family leverages Multimodal Large Language Models (MLLMs) as embedding backbones~\citep{kong2025modality}. Unlike CLIP-based models with separate modality-specific encoders, MLLM-based models employ a \textbf{unified architecture} that processes text, images, and video through the same model. Typically, an MLLM consists of a visual encoder (e.g., a ViT) that converts images or video frames into visual tokens, which are then concatenated with text tokens and jointly processed by a large language model (LLM). The final embedding is extracted via pooling over the output token sequence.

Representative MLLM-based embedding models include UNITE~\citep{kong2025modality}, VLM2Vec-V2~\citep{meng2025vlm2vec}, and LLaVE~\citep{lan2025llave}. These models inherit the strong language understanding and visual generalization capabilities of their LLM backbones, enabling them to handle diverse query formats---textual, composed, and visual---within a single model. This architectural shift from modality-specific encoders to unified MLLM-based encoders reflects a broader trend toward general-purpose video embedding models~\citep{qwen3vlembedding}.

\subsection{Spatiotemporal Sampling}
\label{sec:bc-spatiotemporal}

Before a video can be processed by an embedding model, it must be converted into a discrete sequence of visual inputs through \textbf{spatiotemporal sampling}. Since encoding every frame at full resolution is computationally prohibitive~\citep{hu2022mmsampler}, this preprocessing step critically determines both the available visual evidence and the retrieval quality~\citep{guo2025gve}.

\textbf{Temporal sampling} selects $F$ frames from a video of duration $T$. The most common strategy is uniform sampling, where frames are evenly spaced across the video duration.

\textbf{Spatial sampling} determines the visual resolution or token representation per frame. For ViT-based architectures, this involves converting frames into sequences of visual tokens, where the token count $S$ determines the spatial detail preserved.

The combination of temporal frame count $F$ and spatial token count $S$ per frame defines the total input size $F \times S$ for the embedding model. This decomposition into orthogonal spatiotemporal dimensions allows systems to trade off temporal coverage against spatial detail under a fixed computational budget~\citep{zhu2025dynamic}.

\subsection{Approximate Nearest Neighbor Search}
\label{sec:bc-anns}

Given a query embedding $\mathbf{q} \in \mathbb{R}^d$ and a large-scale database $\mathcal{D} = \{\mathbf{v}_1, \mathbf{v}_2, \ldots, \mathbf{v}_N\}$ of $N$ candidate embeddings, the nearest neighbor search task is to find the $K$ items in $\mathcal{D}$ most similar to $\mathbf{q}$. When $N$ is large (millions to billions), exact search via exhaustive comparison becomes computationally prohibitive, motivating the development of \textbf{Approximate Nearest Neighbor Search (ANNS)} algorithms that achieve sub-linear query time by building index structures for selective exploration of the database.

Three main categories of index structures have been developed for ANNS:

\textbf{Tree-based methods} recursively partition the vector space using hierarchical structures. KD-trees~\citep{friedman1977algorithm} perform well in low-dimensional spaces but suffer from the ``curse of dimensionality'' as dimensionality increases.

\textbf{Hashing-based methods} map high-dimensional vectors to compact binary codes using locality-sensitive hashing (LSH)~\citep{indyk1998approximate}, enabling fast similarity estimation through Hamming distance.

\textbf{Graph-based methods} construct a proximity graph where nodes represent database items and edges connect similar items. During retrieval, a greedy traversal starting from entry points iteratively expands the candidate set by examining neighbors of visited nodes. The Hierarchical Navigable Small World (HNSW) algorithm~\citep{malkov2018efficient} builds a multi-layered graph where upper layers provide sparse long-range connections for coarse navigation and lower layers provide dense local connections for fine-grained search. Graph-based methods have demonstrated superior performance on high-dimensional data and are widely adopted in production systems~\citep{wang2021comprehensive}.

A key property for efficient graph traversal is the \textbf{triangle inequality}: for any three points $a$, $b$, $c$ in a metric space, $d(a, c) \leq d(a, b) + d(b, c)$. When this property holds, greedy graph traversal is guaranteed to follow optimal paths. However, learned neural relevance metrics often violate the triangle inequality~\citep{chen2022approximate}, forcing the exploration of redundant edges and degrading retrieval efficiency. Designing search-compatible metrics that approximately satisfy this property is important for effective graph-based retrieval.

\subsection{Industrial Retrieval Systems}
\label{sec:bc-industrial}

Modern large-scale recommendation systems adopt a \textbf{retrieval-and-ranking architecture}~\citep{cheng2016wide,covington2016deep}. The retrieval stage reduces a billion-scale item catalog to a thousand-scale candidate set using lightweight models, which then serves as input for more sophisticated ranking models. This multi-stage design is essential for balancing effectiveness and efficiency at web scale~\citep{liu2021jizhi}.

The \textbf{two-tower architecture}~\citep{yang2020mixed} is widely adopted for the retrieval stage. It encodes queries and items through separate neural networks into embeddings compared via efficient similarity metrics such as inner product or cosine similarity. This architecture enables pre-computation and indexing of item representations, allowing real-time retrieval through ANNS.

Beyond simple distance metrics, recent works leverage \textbf{DNN-based learnable metrics}~\citep{dehghani2017neural,tay2018latent} to estimate user-item relevance through high-dimensional computation over user and item embeddings, capturing more complex interaction patterns than hand-designed similarity functions~\citep{tan2020fast,chen2022approximate}. Representative systems include YouTube's recommendation engine~\citep{covington2016deep}, Pinterest's Pixie~\citep{eksombatchai2018pixie}, and LinkedIn's CaSMoS~\citep{borisyuk2016casmos}.

\section{Literature Review}
\label{sec:rw-literature}

This section surveys existing work along the video retrieval pipeline, covering the evolution of retrieval benchmarks and embedding models (\S\ref{sec:lr-retrieval}), the challenges of spatiotemporal scaling for long-video retrieval (\S\ref{sec:lr-spatiotemporal}), and the integration of learned retrieval metrics with large-scale graph-based search systems (\S\ref{sec:lr-graph}).

\subsection{Video Retrieval Benchmarks and Embedding Models}
\label{sec:lr-retrieval}

Text-to-video retrieval has progressed from matching coarse phrases to parsing fine-grained spatio-temporal descriptions. Early datasets established the foundational paradigm: MSR-VTT~\citep{xu2016msr} introduced large-scale video-caption pairing with short clips, while DiDeMo~\citep{anne2017localizing} focused on temporal moment localization. Subsequent benchmarks extended the scope in complementary directions: CaReBench~\citep{xu2025carebench} incorporates spatio-temporal grounding annotations, UltraVideo~\citep{ultravideo} targets scene understanding, and CameraBench~\citep{lin2025towards} evaluates camera motion comprehension. Beyond text-to-video retrieval, new retrieval paradigms have emerged, including composed queries with text-image pairs (e.g., CoVR~\citep{ventura2024covr,hummel2024egocvr}) and purely visual queries in egocentric contexts~\citep{liu2021activity}. For long-video retrieval, LoVR~\citep{cai2025lovr} targets retrieval over hour-long videos, while TVR~\citep{lei2020tvr} and ActivityNet Captions~\citep{krishna2017dense} provide fine-grained temporal annotations for partially relevant video retrieval~\citep{dong2022partially}. Despite this progress, evaluation remains fragmented, with models optimized and tested on isolated benchmarks, preventing a holistic understanding of model generalization across heterogeneous retrieval settings.

The development of video embedding models can be traced through two major phases. The first phase extended CLIP~\citep{radford2021learning} to video by introducing temporal modeling: CLIP4Clip~\citep{luo2022clip4clip} pioneered temporal aggregation over frame-level features, followed by InternVideo2~\citep{wang2024internvideo2} and X-CLIP~\citep{gorti2022x}, which developed more sophisticated spatiotemporal encoders. These models, all built on the dual-encoder paradigm, achieved strong results on standard text-to-video benchmarks but were inherently constrained by their modality-specific architecture.

The second phase, more recent, shifted to MLLM-based backbones that process multiple modalities within a unified architecture. Models such as LLaVE~\citep{lan2025llave}, UNITE~\citep{kong2025modality}, and VLM2Vec-V2~\citep{meng2025vlm2vec} demonstrated that leveraging the language understanding and visual generalization of large pretrained models could yield strong embedding performance. Qwen3-VL-Embedding~\citep{qwen3vlembedding} further showed the potential of large-scale supervised training for MLLM-based embeddings. Despite these advances, a systematic evaluation of such models across diverse retrieval tasks and semantic domains has been lacking, leaving their true generalization capabilities unclear.

\subsection{Long-Video Retrieval and Spatiotemporal Scaling}
\label{sec:lr-spatiotemporal}

Early text-video retrieval benchmarks such as MSR-VTT~\citep{xu2016msr} and DiDeMo~\citep{anne2017localizing} consist of short clips where sparse frame sampling suffices to capture the full content, rendering sampling strategy largely inconsequential. As retrieval extends to longer and more informationally dense videos, the choice of spatiotemporal sampling becomes critical~\citep{guo2025gve}. Since encoding every frame at full resolution is computationally prohibitive~\citep{hu2022mmsampler}, models must carefully balance temporal coverage and spatial detail under a fixed computational budget.

Scaling-law studies in this domain have primarily examined model size and training data volume, establishing how performance scales with these factors. A complementary effort~\citep{zhu2025dynamic} establishes power-law scaling with respect to frame count and caption density on short clips, providing insights into the data efficiency of video models. Prior work on efficient video modeling has studied sparse sampling strategies, with VideoMAE~\citep{tong2022videomae} demonstrating that masked autoencoders can learn effective representations from sparsely sampled frames by reconstructing masked portions of the input. However, these studies focus on representation learning rather than retrieval performance, and their findings may not transfer to retrieval settings where the goal is to compare videos against queries in a shared embedding space.

How spatiotemporal sampling interacts with retrieval performance across different query granularities---from coarse thematic descriptions to fine-grained moment localization---remains underexplored, particularly for long-form videos where temporal sparsity poses distinct challenges. Intuitively, coarse-grained retrieval may benefit from dense temporal sampling to capture broad video themes, while fine-grained retrieval may require high spatial resolution to capture subtle visual details. However, this intuition has not been systematically validated, and the optimal sampling strategy may vary depending on the query type, video length, and computational constraints. This gap motivates the spatiotemporal scaling analysis proposed in Chapter~\ref{ch:spatiotemporal-scaling}.

\subsection{Learned Metrics and Graph-based Retrieval}
\label{sec:lr-graph}

As defined in \S\ref{sec:bc-anns}, graph-based ANNS methods such as HNSW~\citep{malkov2018efficient} have become the dominant approach for high-dimensional similarity search, with subsequent works developing K-nearest neighbor graphs~\citep{fu2016efanna}, relative neighborhood graphs~\citep{jayaram2019diskann}, and GPU-optimized constructions~\citep{ootomo2024cagra}. These methods achieve sub-linear query time by constructing proximity graphs that enable greedy traversal toward the query, but their effectiveness depends critically on the properties of the similarity metric used during traversal. In parallel, the retrieval metrics deployed in industrial systems have evolved from simple distance functions such as cosine similarity to learned DNN-based metrics that capture more complex relevance patterns (as discussed in \S\ref{sec:bc-industrial}).

A critical challenge arises at the intersection of these two trends: neural relevance metrics often violate the triangle inequality~\citep{chen2022approximate}, a property that ensures optimal pathfinding in graph traversal (see \S\ref{sec:bc-anns}). This violation forces the exploration of redundant edges and degrades retrieval efficiency, as the greedy traversal may follow paths that appear promising according to the learned metric but do not lead to the true nearest neighbors. Early industrial frameworks such as SL2G~\citep{tan2020fast} first leveraged DNN-based metrics for graph-based ANNS in real-world recommendations, demonstrating that learned metrics could improve retrieval accuracy despite their incompatibility with graph traversal. NANN~\citep{chen2022approximate} extended this by improving the DNN-based metric for better accuracy and adaptivity, but the fundamental tension between metric expressiveness and search compatibility remained unresolved.

The co-design of learned retrieval metrics with graph traversal efficiency---ensuring that the metric not only captures relevance accurately but also supports efficient navigation---remains an open problem, particularly in GPU-accelerated environments serving high-concurrency workloads~\citep{liu2022lion,liu2021jizhi}.
